% GENERATED FILE. Edit canonical lessons, then run npm run generate:books.
% canonical-source-hash: fnv1a64:bb7712f990e22bf6
% canonical-lessons: PT-C456-clareira, PT-C456-riacho, PT-C456-ribeira, PT-C456-arbusto, PT-C456-vadear, PT-R456-revisao-449-456

\chapter{No bosque}
\label{ch:pt-no-bosque}
\hlchaptermodality{456}

\begin{chapteropening}
\textbf{By the end of this chapter:} \emph{I can describe a walk in the woods: a clearing, a brook, a stream, a bush and wading across.}

Complete the last lesson of chapter 456: I can describe a walk in the woods: a clearing, a brook, a stream, a bush and wading across.
\end{chapteropening}

\section[a clareira]{a clareira — the clearing (in a wood)}

\label{lesson:PT-C456-clareira}



\begin{quote}
Before the new one: say the Portuguese for the bank statement, then the Portuguese for the fee.
\end{quote}

\subsection*{You'll want to know: a clareira}
\textbf{a clareira} — \textquotedblleft{}the clearing (in a wood)\textquotedblright{}.

\begin{quote}
Fizemos um piquenique numa clareira.
\end{quote}

\begin{quote}
\emph{We had a picnic in a clearing.}
\end{quote}

Built on \emph{claro}: the bright place among the trees.

\begin{cousinweb}
\emph{clareira} is built on \emph{claro} 'bright, clear' with \emph{-eira} (EtymDB); \emph{claro} is inherited from Latin \emph{clārus}.
\end{cousinweb}

\begin{grammarlens}[title={What you've built}]
One more word for telling, explaining and describing at B1.
\end{grammarlens}

\subsection*{Your turn}
\begin{itemize}
\raggedright
  \item \emph{Say it:} \emph{a clareira}
  \item \emph{Say it:} \emph{a clareira}, once more
  \item \emph{Say it:} a sentence of your own with \emph{a clareira}
  \item \emph{Recall:} say the Portuguese for the bank statement, then the Portuguese for the fee, then say \emph{a clareira} again
\end{itemize}

\begin{tcolorbox}[breakable,colback=teal!4,colframe=teal!35!black,title={Before you move on}]
What does a clareira mean? (\textquotedblleft{}The clearing (in a wood)\textquotedblright{}.) Say it once more.
\end{tcolorbox}
\section[o riacho]{o riacho — the brook}

\label{lesson:PT-C456-riacho}



\begin{quote}
Before the new one: say the Portuguese for the fee, then the Portuguese for the clearing.
\end{quote}

\subsection*{You'll want to know: o riacho}
\textbf{o riacho} — \textquotedblleft{}the brook\textquotedblright{}.

\begin{quote}
Um riacho passa pelo meio do bosque.
\end{quote}

\begin{quote}
\emph{A brook runs through the middle of the wood.}
\end{quote}

Smaller than \emph{um rio}.

\begin{cousinweb}
\emph{riacho} was borrowed from Spanish \emph{riacho} (EtymDB), a small form of \emph{río} 'river', from Latin \emph{rīvus} 'stream' (Lewis \& Short).
\end{cousinweb}

\begin{grammarlens}[title={What you've built}]
One more word for telling, explaining and describing at B1.
\end{grammarlens}

\subsection*{Your turn}
\begin{itemize}
\raggedright
  \item \emph{Say it:} \emph{o riacho}
  \item \emph{Say it:} \emph{o riacho}, once more
  \item \emph{Say it:} a sentence of your own with \emph{o riacho}
  \item \emph{Recall:} say the Portuguese for the fee, then the Portuguese for the clearing, then say \emph{o riacho} again
\end{itemize}

\begin{tcolorbox}[breakable,colback=teal!4,colframe=teal!35!black,title={Before you move on}]
What does o riacho mean? (\textquotedblleft{}The brook\textquotedblright{}.) Say it once more.
\end{tcolorbox}
\section[a ribeira]{a ribeira — the stream; the riverside}

\label{lesson:PT-C456-ribeira}



\begin{quote}
Before the new one: say the Portuguese for the clearing, then the Portuguese for the brook.
\end{quote}

\subsection*{You'll want to know: a ribeira}
\textbf{a ribeira} — \textquotedblleft{}the stream; the riverside\textquotedblright{}.

\begin{quote}
As crianças apanham pedras na ribeira.
\end{quote}

\begin{quote}
\emph{The children pick up stones in the stream.}
\end{quote}

In old town names \emph{a Ribeira} is the riverside quarter.

\begin{cousinweb}
\emph{ribeira} is inherited from Latin \emph{rīpāria} 'of the riverbank', from \emph{rīpa} 'bank' (Lewis \& Short); EtymDB lists Spanish \emph{ribera} and French \emph{rivière} as cousins.
\end{cousinweb}

\begin{grammarlens}[title={What you've built}]
One more word for telling, explaining and describing at B1.
\end{grammarlens}

\subsection*{Your turn}
\begin{itemize}
\raggedright
  \item \emph{Say it:} \emph{a ribeira}
  \item \emph{Say it:} \emph{a ribeira}, once more
  \item \emph{Say it:} a sentence of your own with \emph{a ribeira}
  \item \emph{Recall:} say the Portuguese for the clearing, then the Portuguese for the brook, then say \emph{a ribeira} again
\end{itemize}

\begin{tcolorbox}[breakable,colback=teal!4,colframe=teal!35!black,title={Before you move on}]
What does a ribeira mean? (\textquotedblleft{}The stream; the riverside\textquotedblright{}.) Say it once more.
\end{tcolorbox}
\section[o arbusto]{o arbusto — the bush, the shrub}

\label{lesson:PT-C456-arbusto}



\begin{quote}
Before the new one: say the Portuguese for the brook, then the Portuguese for the stream.
\end{quote}

\subsection*{You'll want to know: o arbusto}
\textbf{o arbusto} — \textquotedblleft{}the bush, the shrub\textquotedblright{}.

\begin{quote}
Um coelho escondeu-se atrás de um arbusto.
\end{quote}

\begin{quote}
\emph{A rabbit hid behind a bush.}
\end{quote}

Smaller than a tree, with many stems.

\begin{cousinweb}
\emph{arbusto} is a learned borrowing from Latin \emph{arbustum} 'place planted with trees' (EtymDB), built on \emph{arbor} 'tree'.
\end{cousinweb}

\begin{grammarlens}[title={What you've built}]
One more word for telling, explaining and describing at B1.
\end{grammarlens}

\subsection*{Your turn}
\begin{itemize}
\raggedright
  \item \emph{Say it:} \emph{o arbusto}
  \item \emph{Say it:} \emph{o arbusto}, once more
  \item \emph{Say it:} a sentence of your own with \emph{o arbusto}
  \item \emph{Recall:} say the Portuguese for the brook, then the Portuguese for the stream, then say \emph{o arbusto} again
\end{itemize}

\begin{tcolorbox}[breakable,colback=teal!4,colframe=teal!35!black,title={Before you move on}]
What does o arbusto mean? (\textquotedblleft{}The bush, the shrub\textquotedblright{}.) Say it once more.
\end{tcolorbox}
\section[vadear]{vadear — to wade across, to ford}

\label{lesson:PT-C456-vadear}



\begin{quote}
Before the new one: say the Portuguese for the stream, then the Portuguese for the bush.
\end{quote}

\subsection*{You'll want to know: vadear}
\textbf{vadear} — \textquotedblleft{}to wade across, to ford\textquotedblright{}.

\begin{quote}
Vadeámos a ribeira com as botas na mão.
\end{quote}

\begin{quote}
\emph{We waded across the stream with our boots in our hands.}
\end{quote}

Built on \emph{o vau}, the ford.

\begin{cousinweb}
\emph{vadear} goes back to Latin \emph{vadum} 'shallow place, ford', the source of \emph{vau} 'ford' (EtymDB).
\end{cousinweb}

\begin{grammarlens}[title={What you've built}]
One more word for telling, explaining and describing at B1.
\end{grammarlens}

\subsection*{Your turn}
\begin{itemize}
\raggedright
  \item \emph{Say it:} \emph{vadear}
  \item \emph{Say it:} \emph{vadear}, once more
  \item \emph{Say it:} a sentence of your own with \emph{vadear}
  \item \emph{Recall:} say the Portuguese for the stream, then the Portuguese for the bush, then say \emph{vadear} again
\end{itemize}

\begin{tcolorbox}[breakable,colback=teal!4,colframe=teal!35!black,title={Before you move on}]
What does vadear mean? (\textquotedblleft{}To wade across, to ford\textquotedblright{}.) Say it once more.
\end{tcolorbox}
\section[(dialogue)]{Revisão dos capítulos 449–456}

\label{lesson:PT-R456-revisao-449-456}



\begin{quote}
Say the Portuguese for the household goods shop and the haberdashery.
\end{quote}

\subsection*{Your turn}
Say these aloud:

\begin{itemize}
\raggedright
  \item the household goods shop, the haberdashery, the perfume shop, the herbalist's shop, the jeweller's shop
  \item the grocer's, the deli, the fruit shop, the wine shop, to slice
  \item the cobbler, the seamstress, the hem, the locksmith, to patch
  \item the barber, to trim, the beautician, the optician, the stain
  \item the till receipt, the bargain, the clearance sale, to haggle, to wrap
  \item the sender, the addressee, the postage, to dispatch, to register
  \item the cash withdrawal, to charge, to debit, the bank statement, the fee
  \item the clearing, the brook, the stream, the bush, to wade across
\end{itemize}

\begin{tcolorbox}[breakable,colback=teal!4,colframe=teal!35!black,title={Before you move on}]
The household goods shop? (\textbf{a drogaria}.) The haberdashery? (\textbf{a retrosaria}.) The perfume shop? (\textbf{a perfumaria}.) The herbalist's shop? (\textbf{a ervanária}.) The jeweller's shop? (\textbf{a ourivesaria}.) The grocer's? (\textbf{a mercearia}.) The deli? (\textbf{a charcutaria}.) The fruit shop? (\textbf{a frutaria}.) The wine shop? (\textbf{a garrafeira}.) To slice? (\textbf{fatiar}.) The cobbler? (\textbf{o sapateiro}.) The seamstress? (\textbf{a costureira}.) The hem? (\textbf{a bainha}.) The locksmith? (\textbf{o serralheiro}.) To patch? (\textbf{remendar}.) The barber? (\textbf{o barbeiro}.) To trim? (\textbf{aparar}.) The beautician? (\textbf{a esteticista}.) The optician? (\textbf{o oculista}.) The stain? (\textbf{a nódoa}.) The till receipt? (\textbf{o talão}.) The bargain? (\textbf{a pechincha}.) The clearance sale? (\textbf{a liquidação}.) To haggle? (\textbf{regatear}.) To wrap? (\textbf{embrulhar}.) The sender? (\textbf{o remetente}.) The addressee? (\textbf{o destinatário}.) The postage? (\textbf{os portes}.) To dispatch? (\textbf{expedir}.) To register? (\textbf{registar}.) The cash withdrawal? (\textbf{o levantamento}.) To charge? (\textbf{cobrar}.) To debit? (\textbf{debitar}.) The bank statement? (\textbf{o extrato}.) The fee? (\textbf{a comissão}.) The clearing? (\textbf{a clareira}.) The brook? (\textbf{o riacho}.) The stream? (\textbf{a ribeira}.) The bush? (\textbf{o arbusto}.) To wade across? (\textbf{vadear}.)
\end{tcolorbox}
